\documentclass{article}

\PassOptionsToPackage{numbers}{natbib}
\usepackage[final]{neurips_2019}

\usepackage[utf8]{inputenc}
\usepackage[T1]{fontenc}
\usepackage{hyperref}
\usepackage{url}
\usepackage{booktabs}
\usepackage{amsfonts}
\usepackage{amsmath}
\usepackage{microtype}
\usepackage{graphicx}
\usepackage{xcolor}
\usepackage{xparse}

\newif\ifshowinstructions
\showinstructionsfalse

\newcommand{\note}[1]{\ifshowinstructions#1\fi}
\NewDocumentEnvironment{instructions}{+b}
  {\ifshowinstructions #1\par\fi}
  {}

\NewDocumentCommand{\projecttitle}{m m m O{} O{}}{%
  \title{%
    #1\\[0.7em]
    {\normalsize\normalfont
      Columbia COMS 4705 Project Milestone
    }\\[0.8em]
    \parbox{0.95\textwidth}{%
      \normalfont\small\centering

      \textbf{Keywords:} \textit{#2}
      \par\vspace{0.4em}

      \textbf{TA Mentor:} #3%

      \if\relax\detokenize{#4}\relax\else
        \par\vspace{0.4em}
        \textbf{External Collaborators:} #4%
      \fi

      \if\relax\detokenize{#5}\relax\else
        \par\vspace{0.4em}
        \textbf{Sharing Project:} #5%
      \fi
      \par
    }%
  }%
}

\projecttitle
  {What Does a Wish Leave Out? Requirement Recovery in Wish-Driven Code Agents}
  {wish-driven programming, code agents, requirement elicitation, underspecification}
  {Noah Foster}
  []
  []

\author{
  Yanhao Bai \quad Kaitlyn Bassford \quad Wen Tao Koh \\
  Department of Computer Science, Columbia University \\
  \texttt{\{yb2630, kmb2297, wk2439\}@columbia.edu}
}

\begin{document}

\maketitle

\begin{abstract}
Wish-driven programming---stating an intent and letting an agent build the artifact---fails
quietly: the program runs, passes the tests the user remembered to mention, and violates a
requirement nobody wrote down. We degrade benchmark specifications into wishes by deleting
fields mechanically, and train a small \emph{silence detector} that predicts which requirement
categories a wish leaves implicit and uses that ranking to choose which clarifying question to
ask. On 100 development tasks with Qwen3-8B, targeted questioning reaches 53.7\% against 44.3\%
for generic questioning at the same budget and 42.7\% when the same answers are withheld, while
a harness that may read neighbouring code reaches 51.7\%.
\end{abstract}

\section{Approach}

\paragraph{The task, and what actually differs across the three datasets.}
Each task is one Python function: the agent sees a wish of one or two sentences, a stub, and
three visible tests that print observed outputs but never expected ones; it must pass a hidden
suite taken from EvalPlus \citep{liu2023evalplus} or BigCodeBench. The harness is deliberately
thin---a ReAct loop with four tools (write the file, run the visible tests, read the last
error, ask the stakeholder), at most twelve turns, and \emph{no repository}. A codebase would
let the agent recover the deleted requirements from neighbouring code and destroy the
manipulation, so rather than assert that, condition H4 supplies five sibling functions and
measures it. Contemporary harnesses also revise their own state online \citep{xiao2026reharness};
we hold state fixed so that only the prompt varies. The three datasets are not interchangeable. BigCodeBench contributes 591 tasks, every one of
them constraining its arguments and 25.7\% also specifying error behaviour; MBPP+ contributes
377, with those fields present in 8.2\% and 0.3\%, and HumanEval+ 132, at 14.4\% and 3.8\%.

\paragraph{Levels are built by deletion, and it is cheap.}
Each docstring is parsed into five fields---intent, signature and types, examples, constraints
and edge cases, error behaviour---and levels delete fields: L0 keeps all, L1 drops examples and
type detail, L2 also drops constraints and error behaviour, L3 keeps the intent alone (the whole
prose block, not its first sentence, after R2 showed that one sentence often names only the
input). Hidden tests are identical across levels, so an L0--L3 difference is information loss, and
the deleted fields are the detector's labels. Building the corpus costs no model calls: 1{,}682 rows, 1{,}651 parsed (31 unparsable), 1{,}549 validated by running their visible and hidden suites (666\,s, eight cores), 1{,}100 kept; the 133 exclusions are 31 parse, 86 visible and 16 hidden. Eleven minutes and \$0 per rebuild is why we could rebuild every level after a parsing bug.

\paragraph{How the detector changes the generated program.}
The detector never touches model weights; it changes what the generator is conditioned on. It
predicts which fields L0 carries and L3 drops, and that ranking chooses which of five question
templates is asked, following the stakeholder protocol of \citep{mu2023clarifygpt}. The
stakeholder's answers are compiled into a \emph{Requirements} block that
(i) conditions the code prompt, (ii) is reused as acceptance criteria in one
self-check-and-repair pass, and (iii) supplies a defensive-coding instruction for any category left
unanswered when the budget runs out. The measured contrast is the mechanism:
targeted questions reach 53.7\%, generic questions at the same budget 44.3\%, and the same
questions with the answers withheld 42.7\%. Targeting is worth $+9.3$ points over generic and
$+11.0$ over withholding, but only $+1.7$ over asking nothing---so the detector's contribution
is to make elicitation worth its cost, not to beat a silent agent.

\paragraph{Conditions and baselines.}
H0 is the plain loop; H1 writes acceptance tests first; H2 asks generic questions; H3 asks
detector-targeted questions; H4 exposes five sibling functions with doctests; H3$^-$ asks H3's questions and
withholds the answers. B1 gives the full specification with no dialogue; B2 draws five
single-shot samples (51.3\% pass@1, 54.7\% over five samples); B3 is a one-call, no-tools baseline
at 31.0\%. H5, not yet run, evolves the harness offline on 200 development wishes with RRSI's
component-wise edits and regularized selection of the RRSI line of work \citep{xia2026rrsi}. The
loop conditions (B1, H0--H4, H3$^-$) share tasks, tools and hidden tests; B2 and B3 are
dialogue-free baselines.

\section{Experiments}

\paragraph{Data, metrics, implementation.}
From 1{,}100 parsed tasks, 500 train the detector, 200 are reserved for harness evolution, 300
held-out tasks are untouched, and the 100-task development split carries every number below. \emph{Recovery} is the share of the L0--L3 gap a condition closes; failures are
bucketed by the exception the runner raises. The generator is Qwen3-8B-AWQ on vLLM 0.34,
temperature 0.7, top-$p$ 0.95, twelve turns, three samples per task, seed 20261121; the run took
2.9 GPU-hours on one A100, and the detector is a linear head on frozen Qwen3-Embedding-0.6B
features (five labels, 600 full-batch AdamW steps, calibrated out of fold).

\begin{center}
  \small
  \begin{tabular}{@{}llrrr@{}}
    \toprule
    Condition & Description & pass@1 & Recovery & Questions \\
    \midrule
    B1        & full specification, no dialogue & 62.3 & 1.00 & 0 \\
    H0        & ReAct, no questions             & 52.0 & 0.00 & 0 \\
    H1        & acceptance tests first          & 28.7 & $-2.26$ & 0 \\
    H2        & generic clarification           & 44.3 & $-0.74$ & 2.94 \\
    H3        & detector-targeted clarification & \textbf{53.7} & 0.16 & 2.48 \\
    H4        & context-rich (sibling functions)& 51.7 & $-0.03$ & 0 \\
    H3$^-$    & H3's questions, answers withheld& 42.7 & $-0.90$ & 2.48 \\
    \bottomrule
  \end{tabular}
  \vspace{2pt}

  {\footnotesize\raggedright\textbf{Table 2.} R2b pilot, 100 development tasks, 3 samples,
  single seed; held-out tasks used: 0.\par}
\end{center}

\paragraph{Results.}
The specification--wish gap is 10.3 points and no condition closes it: H3, the best dialogue
condition, recovers 16\%. Task-level bootstrap over all three samples (2{,}000 resamples, seed
20261121) gives the contrasts: targeting beats generic questions by $+9.3$ pp $[+1.7,\,+17.0]$;
withholding H3's answers costs $11.0$ pp $[-18.7,\,-2.7]$; H3 is not distinguishable from silence
($+1.7$ pp, $[-5.7,\,+9.0]$); acceptance tests first costs $23.3$ pp $[-31.3,\,-15.7]$; generic
questions cost 7.7 pp $[-14.3,\,-1.0]$. The context-rich harness neither helps nor
hurts (51.7)---browsing is not a substitute for eliciting. Failures are mostly wrong outputs
(\texttt{AssertionError} 858); the detector reaches macro-F1 0.792 (TF-IDF 0.789) and AUROC 0.744
on error behaviour, where the label is rare (85 of 500 training tasks).

Code: we wrote the level builder, harness, detector and evaluation; the hidden tests are EvalPlus's and BigCodeBench's, and the evolution search follows RRSI.

\section{Future Work}

Since the proposal, three things changed: the harness is now specified per dataset, levels are
built mechanically (11 minutes, \$0), and the detector's effect on generation is measured rather
than asserted. Three things remain: the 300 held-out tasks and a second seed; H5, to test whether
an evolved harness transfers or overfits; and a refit on the corrected labels---two of the five
are near-constant, so H3's $+9.3$ should be re-measured against a three-label detector.

\bibliographystyle{unsrt}
\bibliography{references}

\end{document}
